\Titular%
{miscelanea}%
{Música, frases e máis}%
{}%
{}%

\vspace*{3mm}

\subsection*{\textcolor{Resalte}{Praceres descoñecidos: Unknown Pleasures de Joy Division \\ } {\normalfont \itshape   Carlos Carballeira}}%

\begin{multicols}{2}

Aínda que este número de \textit{Momentum} chegará días despois ás vosas mans, estou escribindo estas liñas en vésperas do día máis triste do ano: o \textit{Blue Monday}. Así quedou bautizado fará uns vinte anos o terceiro luns do mes de xaneiro, logo de que un profesor dun centro asociado á Universidade de Cardiff, chamado Cliff Arnall, calculase (desenvolveu unha fórmula para tal fin...) que a combinación do rigor climático do inverno, a costa de Xaneiro, e a depresión post-Nadal, facían de esa data unha xornada especialmente deprimente.

\begin{figure}[H] 
    \centering
    \includegraphics[width=0.7\linewidth]{./revistas/005/imaxes/u_pl.jpg}
    \caption*{Portada orixinal de \textit{Unknown Pleasures}, álbum de debut de \textit{Joy Division} (1979). A imaxe, que mostra unha secuencia de pulsos de radiofrecuencia emitidos polo primeiro púlsar descuberto (o PSR B1919+21), é unha das máis icónicas da historia da música.}
\end{figure}

Antes da ocorrencia de Arnall, \textit{Blue Monday} foi o título da canción coa que a banda de Manchester \textit{New Order} revolucionou a música electrónica e de baile no ano 1983. O uso pioneiro que fixo \textit{New Order} dos sintetizadores e da programación rítmica na súa discografía en xeral e nese tema en particular, foi un fito clave na transición entre o \textit{post-punk} e xéneros como o \textit{techno}, o \textit{house} e o \textit{synth-pop}. A listaxe de artistas que recoñeceron explicitamente a súa influencia inclúe, entre outros moitos, a \textit{Moby}, \textit{The Chemical Brothers} ou \textit{The Killers}. De feito, estes últimos adoptaron o seu nome dunha banda ficticia que aparece no vídeo-clip do tema \textit{Crystal} de \textit{New Order} (2001). Moi preto de nos, o grupo \textit{indie} de Boiro \textit{Triangulo de Amor Bizarro} leva tamén o nome de outro dos grandes éxitos de \textit{New Order}. 

\textit{New Order} emerxeu das cinzas doutra banda que, aínda que efimera, foi tan ou máis influínte na música posterior: \textit{Joy Division}. O seu álbum de debut, \textit{Unknown Pleasures} (1979), definiu unha nova estética musical que transitaba dende a exuberancia e os excesos do \textit{punk} ata un estilo moito máis contido, con bases rítmicas melódicas, letras instrospectivas e unha producción moi traballada que foi definido como \textit{post-punk}. Pero ademais, ese álbum tamén ten conexión coa (astro-)física non por unha cancion, se non pola súa portada: Stephen Morris, batería da banda, atopou unha imaxe hipnótica na \textit{Cambridge Encyclopedia of Astronomy} e amosoulla a Peter Saville, deseñador da discográfica \textit{Factory Records} (recoméndovos a película \textit{24 hour party people} sobre a historia desa compañía). Invertendo as cores do diagrama orixinal, Saville creou unha imaxe minimalista, sen título nin nome da banda, de liñas brancas sobre fondo negro que xeraba misterio arredor do disco. O misterio permaneceu durante décadas, e non foi ata 2011 cando se descubriu que tal imaxe fora orixinalmente publicada en \textit{Scientific American} no ano 1971 polo astrónomo Harold D. Craft, quen a creou como parte da súa tese de doutoramento na que analizaba os pulsos de radiofrecuencia do primeiro púlsar detectado na historia (o PSR B1919+21). Os púlsares xiran sobre si mesmos a gran velocidade emitindo ondas de radio, como un faro no espazo, e cando a terra interponse na súa traxectoria detectase un pulso nese rango do espectro. A imaxe creada por Craft era a representación de varios pulsos consecutivos amoreados en vertical.    

A historia de \textit{Joy Division} foi curta coma un pulso, xa que a prematura morte de Ian Curtis (cantante e letrista) desencadeou, despois do disco póstumo \textit{Closer}, a súa desaparición e reinvención en forma de \textit{New Order}. Con  todo, \textit{Unknown Pleasures} tivo unha influencia decisiva en xéneros como o \textit{indie}, o gótico e o rock alternativo en xeral, e a súa portada converteuse nunha icona popular moito antes de que se coñecera a súa orixe. De feito, non é raro ver xente vestindo camisetas con esa imaxe incluso pola nosa facultade. Ata Primark chegou a comercializalas, como tamén fixo con portadas de discos de bandas como \textit{Nirvana} ou \textit{AC/DC}. Velaquí a súa orixe... e tamén a do nome de este recuncho.


\end{multicols}

